\begin{description}
  \item[\texttt{CMS-first} архитектура] Подход, при котором \texttt{CMS} выступает центром управления контентом, а публичная витрина реализуется как отдельное приложение.
  \item[\texttt{Dynamic Zone}] Механизм \texttt{Strapi}, позволяющий собирать страницу из переиспользуемых блоков разного типа.
  \item[\texttt{Headless CMS}] Система управления контентом, отделяющая хранение и редактирование данных от их публичного отображения.
  \item[Контент-модель] Набор сущностей, полей и связей, определяющий структуру данных в CMS.
  \item[Локализованный маршрут] Публичный URL, зависящий от выбранной локали, например \texttt{/ru/...} или \texttt{/en/...}.
  \item[Пайплайн сборки] Последовательность автоматических шагов проверки, сборки и публикации проекта.
  \item[Публикационный контур] Согласованная цепочка действий и компонентов, через которую опубликованный в CMS контент попадает на публичную витрину.
  \item[Публичная витрина] Frontend-часть системы, которая получает данные из CMS и публикует их в виде пользовательских веб-страниц.
  \item[Режим предварительного просмотра] Механизм безопасного просмотра черновой версии материала до его публикации.
  \item[Роутинг] Механизм сопоставления URL-адресов с обработчиками или страницами приложения.
  \item[Серверный рендеринг] Формирование HTML на сервере перед отправкой результата клиенту.
  \item[\texttt{Slug}] Человекочитаемый идентификатор страницы или записи, используемый в составе URL-адреса.
  \item[Типизация контента] Формальное описание структуры данных для валидации и безопасного обмена.
  \item[\texttt{Webhook}] HTTP-вызов, автоматически отправляемый одной системой другой при наступлении определенного события.
\end{description}
